% GENERATED FILE. Edit canonical lessons, then run npm run generate:books.
% canonical-source-hash: fnv1a64:08099fceaa7d0f8b
% canonical-lessons: BN-C539-bighno, BN-C539-nachhorbanda, BN-C539-moriya, BN-C539-hal-chhara, BN-C539-jorano

\chapter{Facing Obstacles}
\label{ch:bn-facing-obstacles}
\hlchaptermodality{539}

\begin{chapteropening}
\textbf{By the end of this chapter:} \emph{I can talk about obstacles, getting tangled up and not giving up.}

Complete the last lesson of chapter 539: I can talk about obstacles, getting tangled up and not giving up.
\end{chapteropening}

\section[\texorpdfstring{\bn{বিঘ্ন} (bighnô)}{বিঘ্ন (bighnô)}]{\bn{বিঘ্ন} (bighnô) — an obstacle; a hindrance}

\label{lesson:BN-C539-bighno}



\begin{quote}
Before the new one: say the Bengali for imminent, then the Bengali for for emergencies.
\end{quote}

\subsection*{You'll want to know: \bn{বিঘ্ন}}
\textbf{\bn{বিঘ্ন}} — \emph{bighnô} — \textquotedblleft{}an obstacle; a hindrance\textquotedblright{}.

\begin{quote}
\bn{বৃষ্টি} \bn{হলে} \bn{আমাদের} \bn{যাত্রায়} \bn{বিঘ্ন} \bn{ঘটবে।}
\end{quote}

\begin{quote}
\emph{Brishṭi hôle āmāder jātrāy bighnô ghôṭbe.}
\end{quote}

\begin{quote}
If it rains, our journey will be held up.
\end{quote}

With \bn{ঘটা}: \bn{বিঘ্ন} \bn{ঘটা}, to be disrupted; \bn{বিঘ্ন} \bn{ঘটানো}, to disrupt.

\begin{cousinweb}
\textbf{\bn{বিঘ্ন}} (\emph{bighnô}) is a tatsama, Sanskrit \emph{vighna}, \textquotedblleft{}a breaker\textquotedblright{}, and so \textquotedblleft{}an obstacle\textquotedblright{} (Monier-Williams).
\end{cousinweb}

\begin{grammarlens}[title={What you've built}]
One more word for saying what happens if.
\end{grammarlens}

\subsection*{Your turn}
\begin{itemize}
\raggedright
  \item \emph{Say it:} \emph{bighnô}
  \item \emph{Say it:} \emph{bighnô}, once more
  \item \emph{Say it:} \emph{Brishṭi hôle āmāder jātrāy bighnô ghôṭbe.} — If it rains, our journey will be held up.
  \item \emph{Recall:} say the Bengali for imminent, then the Bengali for for emergencies, then say \emph{bighnô} again
\end{itemize}

\begin{tcolorbox}[breakable,colback=teal!4,colframe=teal!35!black,title={Before you move on}]
What does \bn{বিঘ্ন} mean? (\textquotedblleft{}An obstacle; a hindrance\textquotedblright{}.) Say it once more.
\end{tcolorbox}
\section[\texorpdfstring{\bn{নাছোড়বান্দা} (nāchhoṛbāndā)}{নাছোড়বান্দা (nāchhoṛbāndā)}]{\bn{নাছোড়বান্দা} (nāchhoṛbāndā) — persistent; never giving up}

\label{lesson:BN-C539-nachhorbanda}



\begin{quote}
Before the new one: say the Bengali for for emergencies, then the Bengali for a hindrance.
\end{quote}

\subsection*{You'll want to know: \bn{নাছোড়বান্দা}}
\textbf{\bn{নাছোড়বান্দা}} — \emph{nāchhoṛbāndā} — \textquotedblleft{}persistent; never giving up\textquotedblright{}.

\begin{quote}
\bn{ছেলেটা} \bn{নাছোড়বান্দা}, \bn{শেষে} \bn{মা} \bn{রাজি} \bn{হলেন।}
\end{quote}

\begin{quote}
\emph{Chheleṭā nāchhoṛbāndā, sheshe mā rāji hôlen.}
\end{quote}

\begin{quote}
The boy wouldn't give up, and in the end his mother agreed.
\end{quote}

\begin{cousinweb}
\textbf{\bn{নাছোড়বান্দা}} (\emph{nāchhoṛbāndā}) is an expressive compound: \bn{না}, \textquotedblleft{}not\textquotedblright{}, \bn{ছোড়}, \textquotedblleft{}letting go\textquotedblright{} (from the family of \bn{ছাড়া}, CDIAL 4998), and \bn{বান্দা}, usually traced to Persian \emph{banda}, \textquotedblleft{}a servant\textquotedblright{}: one who will not let go.
\end{cousinweb}

\begin{grammarlens}[title={What you've built}]
One more word for saying what happens if.
\end{grammarlens}

\subsection*{Your turn}
\begin{itemize}
\raggedright
  \item \emph{Say it:} \emph{nāchhoṛbāndā}
  \item \emph{Say it:} \emph{nāchhoṛbāndā}, once more
  \item \emph{Say it:} \emph{Chheleṭā nāchhoṛbāndā, sheshe mā rāji hôlen.} — The boy wouldn't give up, and in the end his mother agreed.
  \item \emph{Recall:} say the Bengali for for emergencies, then the Bengali for a hindrance, then say \emph{nāchhoṛbāndā} again
\end{itemize}

\begin{tcolorbox}[breakable,colback=teal!4,colframe=teal!35!black,title={Before you move on}]
What does \bn{নাছোড়বান্দা} mean? (\textquotedblleft{}Persistent; never giving up\textquotedblright{}.) Say it once more.
\end{tcolorbox}
\section[\texorpdfstring{\bn{মরিয়া} (môriyā)}{মরিয়া (môriyā)}]{\bn{মরিয়া} (môriyā) — desperate; ready to risk anything}

\label{lesson:BN-C539-moriya}



\begin{quote}
Before the new one: say the Bengali for a hindrance, then the Bengali for persistent.
\end{quote}

\subsection*{You'll want to know: \bn{মরিয়া}}
\textbf{\bn{মরিয়া}} — \emph{môriyā} — \textquotedblleft{}desperate; ready to risk anything\textquotedblright{}.

\begin{quote}
\bn{বন্যার} \bn{জল} \bn{বাড়ছিল}, \bn{তাই} \bn{লোকজন} \bn{মরিয়া} \bn{হল।}
\end{quote}

\begin{quote}
\emph{Bônyār jôl bāṛchhilô, tāi lokjôn môriyā hôlô.}
\end{quote}

\begin{quote}
The flood water was rising, so people grew desperate.
\end{quote}

\begin{cousinweb}
\textbf{\bn{মরিয়া}} (\emph{môriyā}) is probably built on \bn{মরা}, \textquotedblleft{}to die\textquotedblright{} (CDIAL 9871): ready to face death itself.
\end{cousinweb}

\begin{grammarlens}[title={What you've built}]
One more word for saying what happens if.
\end{grammarlens}

\subsection*{Your turn}
\begin{itemize}
\raggedright
  \item \emph{Say it:} \emph{môriyā}
  \item \emph{Say it:} \emph{môriyā}, once more
  \item \emph{Say it:} \emph{Bônyār jôl bāṛchhilô, tāi lokjôn môriyā hôlô.} — The flood water was rising, so people grew desperate.
  \item \emph{Recall:} say the Bengali for a hindrance, then the Bengali for persistent, then say \emph{môriyā} again
\end{itemize}

\begin{tcolorbox}[breakable,colback=teal!4,colframe=teal!35!black,title={Before you move on}]
What does \bn{মরিয়া} mean? (\textquotedblleft{}Desperate; ready to risk anything\textquotedblright{}.) Say it once more.
\end{tcolorbox}
\section[\texorpdfstring{\bn{হাল} \bn{ছাড়া} (hāl chhāṛā)}{হাল ছাড়া (hāl chhāṛā)}]{\bn{হাল} \bn{ছাড়া} (hāl chhāṛā) — to give up hope; to stop trying}

\label{lesson:BN-C539-hal-chhara}



\begin{quote}
Before the new one: say the Bengali for persistent, then the Bengali for desperate.
\end{quote}

\subsection*{You'll want to know: \bn{হাল} \bn{ছাড়া}}
\textbf{\bn{হাল} \bn{ছাড়া}} — \emph{hāl chhāṛā} — \textquotedblleft{}to give up hope; to stop trying\textquotedblright{}.

\begin{quote}
\bn{হাল} \bn{না} \bn{ছাড়লে} \bn{তোমার} \bn{দল} \bn{জিতবে।}
\end{quote}

\begin{quote}
\emph{Hāl nā chhāṛle tomār dôl jitbe.}
\end{quote}

\begin{quote}
If your team doesn't give up, it will win.
\end{quote}

Past (he/she …): \bn{তিনি} \bn{হাল} \bn{ছাড়লেন} \bn{না}, \emph{tini hāl chhāṛlen nā}.

\begin{cousinweb}
\textbf{\bn{হাল} \bn{ছাড়া}} (\emph{hāl chhāṛā}) joins \bn{হাল}, \textquotedblleft{}the helm, the rudder\textquotedblright{}, to \bn{ছাড়া}, \textquotedblleft{}to let go\textquotedblright{}: to let go of the helm. Where \bn{হাল} in this sense comes from is not settled.
\end{cousinweb}

\begin{grammarlens}[title={What you've built}]
One more word for saying what happens if.
\end{grammarlens}

\subsection*{Your turn}
\begin{itemize}
\raggedright
  \item \emph{Say it:} \emph{hāl chhāṛā}
  \item \emph{Say it:} \emph{hāl chhāṛā}, once more
  \item \emph{Say it:} \emph{Hāl nā chhāṛle tomār dôl jitbe.} — If your team doesn't give up, it will win.
  \item \emph{Recall:} say the Bengali for persistent, then the Bengali for desperate, then say \emph{hāl chhāṛā} again
\end{itemize}

\begin{tcolorbox}[breakable,colback=teal!4,colframe=teal!35!black,title={Before you move on}]
What does \bn{হাল} \bn{ছাড়া} mean? (\textquotedblleft{}To give up hope; to stop trying\textquotedblright{}.) Say it once more.
\end{tcolorbox}
\section[\texorpdfstring{\bn{জড়ানো} (jôṛāno)}{জড়ানো (jôṛāno)}]{\bn{জড়ানো} (jôṛāno) — to get tangled up; to wrap round}

\label{lesson:BN-C539-jorano}



\begin{quote}
Before the new one: say the Bengali for desperate, then the Bengali for to give up hope.
\end{quote}

\subsection*{You'll want to know: \bn{জড়ানো}}
\textbf{\bn{জড়ানো}} — \emph{jôṛāno} — \textquotedblleft{}to get tangled up; to wrap round\textquotedblright{}.

\begin{quote}
\bn{ঘুড়ির} \bn{সুতো} \bn{গাছের} \bn{ডালে} \bn{জড়াল।}
\end{quote}

\begin{quote}
\emph{Ghuṛir suto gāchher ḍāle jôṛālô.}
\end{quote}

\begin{quote}
The kite string got tangled in the branches of the tree.
\end{quote}

Past (he/she …): \bn{সেই} \bn{দড়ি} \bn{চাকায়} \bn{জড়াল}, \emph{sei dôṛi chākāy jôṛālô}.

\begin{cousinweb}
\textbf{\bn{জড়ানো}} (\emph{jôṛāno}) comes from a reconstructed \emph{jaḍati}, \textquotedblleft{}joins, sets\textquotedblright{}, which Turner took from Dravidian; CDIAL lists Bengali \emph{jaṛāna}, \textquotedblleft{}to wrap round, to entangle\textquotedblright{} (CDIAL 5091).
\end{cousinweb}

\begin{grammarlens}[title={What you've built}]
One more word for saying what happens if.
\end{grammarlens}

\subsection*{Your turn}
\begin{itemize}
\raggedright
  \item \emph{Say it:} \emph{jôṛāno}
  \item \emph{Say it:} \emph{jôṛāno}, once more
  \item \emph{Say it:} \emph{Ghuṛir suto gāchher ḍāle jôṛālô.} — The kite string got tangled in the branches of the tree.
  \item \emph{Recall:} say the Bengali for desperate, then the Bengali for to give up hope, then say \emph{jôṛāno} again
\end{itemize}

\begin{tcolorbox}[breakable,colback=teal!4,colframe=teal!35!black,title={Before you move on}]
What does \bn{জড়ানো} mean? (\textquotedblleft{}To get tangled up; to wrap round\textquotedblright{}.) Say it once more.
\end{tcolorbox}
